\documentclass[tikz,border=5pt]{standalone}
\usepackage[UTF8]{ctex}
\usepackage{amsmath,amssymb}
\usetikzlibrary{arrows.meta,calc,angles,quotes}
\definecolor{ink}{HTML}{243447}
\definecolor{blue}{HTML}{2765A2}
\definecolor{green}{HTML}{3F7D5B}
\definecolor{red}{HTML}{B74739}
\tikzset{body/.style={draw=ink,thick},vec/.style={draw=blue,thick,-{Latex}},
 guide/.style={gray!60,dashed},dim/.style={ink,{Latex}-{Latex}},line cap=round,line join=round}
\begin{document}

% 题图：行进状态与一圈基膜在侧面展开图中的对应关系。
\begin{tikzpicture}
  \foreach \s/\tail/\head/\lift in {0/0/1/0,1/0/1.8/0,2/0/1.8/1,3/.8/1.8/0}{
    \begin{scope}[xshift=\s*2.7cm]
      \draw[gray] (-.2,-.25)--(2.1,-.25);
      \draw[body,fill=blue!5] (\tail,0) rectangle (\head,.35);
      \draw[body] (\head,0)--(\head,-.25);
      \ifnum\lift=1 \draw[body] (\tail,0)--(\tail-.08,-.1);
      \else \draw[body] (\tail,0)--(\tail,-.25);\fi
      \ifnum\s<2 \fill[red] (\tail,-.25) circle(1.8pt);
      \else \fill[red] (\head,-.25) circle(1.8pt);\fi
      \node[below] at (.9,-.45) {状态\the\numexpr\s+1\relax};
    \end{scope}
  }
  \draw[vec] (8.6,.8)--(10.2,.8) node[midway,above]{前进};
  \node[below] at (.1,-.85) {尾足};\node[below] at (1,-.85) {头足};
  \node[red] at (6.3,-1.55) {红点表示该阶段固定的足};
  \foreach \k/\h in {0/.8,1/1.7,2/2.5}{
    \pgfmathsetmacro{\per}{sqrt(3.2*3.2-\h*\h)}
    \begin{scope}[xshift=\k*3.5cm,yshift=-4.1cm]
      \draw[body,fill=blue!3] (0,0) rectangle (\h*.65,\per*.65);
      \draw[green,very thick] (0,0)--(\h*.65,\per*.65);
      \draw[dim] (-.25,0)--(-.25,\per*.65) node[midway,left]{$P$};
      \draw[dim] (0,-.3)--(\h*.65,-.3) node[midway,below]{$h$};
      \node[green,above,sloped] at (\h*.325,\per*.325) {$L$};
      \pgfmathsetmacro{\ang}{atan(\per/\h)}
      \draw[blue] (0,.32) arc[start angle=90,end angle=\ang,radius=.32];
      \node[blue] at (.18,.42) {$\theta$};
      \pgfmathsetmacro{\rr}{\per/(2*pi)*.65}
      \begin{scope}[yshift=-1.15cm]
        \fill[blue!5] (0,-\rr)--(\h*.65,-\rr)--(\h*.65,\rr)--(0,\rr)--cycle;
        \draw[body] (0,-\rr)--(\h*.65,-\rr) (0,\rr)--(\h*.65,\rr);
        \draw[body] (0,0) ellipse[x radius=.13,y radius=\rr];
        \draw[body] (\h*.65,0) ellipse[x radius=.13,y radius=\rr];
      \end{scope}
    \end{scope}
  }
\end{tikzpicture}

% 展开关系：只有侧面展开，不加上下底。
\begin{tikzpicture}
  \begin{scope}
    \fill[blue!5] (-.9,0) rectangle (.9,2.2);
    \draw[body] (0,0) ellipse[x radius=.9,y radius=.24];
    \draw[body] (0,2.2) ellipse[x radius=.9,y radius=.24];
    \draw[body] (-.9,0)--(-.9,2.2) (.9,0)--(.9,2.2);
    \draw[green,very thick,domain=0:1,samples=100]
      plot({.9*cos(360*\x)},{2.2*\x+.24*sin(360*\x)});
    \draw[dim] (-1.3,0)--(-1.3,2.2) node[midway,left]{$h$};
    \draw[blue,-{Latex}] (0,2.2)--(.9,2.2) node[midway,above]{$r$};
    \node[green] at (.4,1.05) {$L$};
  \end{scope}
  \draw[vec] (1.5,1.1)--(2.4,1.1) node[midway,above]{展开旋转};
  \begin{scope}[xshift=3.1cm]
    \draw[body,fill=blue!3] (0,0) rectangle (2.2,3);
    \draw[green,very thick] (0,0)--(2.2,3) node[midway,above,sloped]{$L$};
    \draw[dim] (0,-.35)--(2.2,-.35) node[midway,below]{$h$};
    \draw[dim] (2.55,0)--(2.55,3) node[midway,right]{$P=2\pi r$};
    \draw[blue] (0,.5) arc[start angle=90,end angle=53.75,radius=.5];
    \node[blue] at (.2,.66) {$\theta$};
    \draw[ink] (1.95,0)--(1.95,.25)--(2.2,.25);
  \end{scope}
  \node at (2.4,-1) {$L^2=h^2+P^2,\qquad \sin\theta=h/L$};
\end{tikzpicture}

% 无量纲体积包络与固定体积的两个交点。
\begin{tikzpicture}[x=8cm,y=7cm]
  \draw[vec] (0,0)--(1.1,0) node[right]{$x=h/L$};
  \draw[vec] (0,0)--(0,.47) node[above]{$f(x)=x(1-x^2)$};
  \draw[blue,very thick,domain=0:1,samples=100] plot(\x,{\x*(1-\x*\x)});
  \draw[red] (0,.146262)--(1.02,.146262);
  \node[red,left] at (0,.146262) {$0.38\cdot\frac{2}{3\sqrt3}$};
  \foreach \a/\lab in {.14961087/{x_{\min}},.57735027/{1/\sqrt3},.91676526/{x_{\max}}}{
    \draw[guide] (\a,0)--(\a,{\a*(1-\a*\a)});
    \node[below] at (\a,0) {$\lab$};
  }
  \fill[green] (.14961087,.146262) circle(2pt) (.91676526,.146262) circle(2pt);
  \draw[green,very thick] (.14961087,.03)--(.91676526,.03);
  \node[green,above] at (.53,.03) {固定体积允许的长度区间};
\end{tikzpicture}

% 等周关系：同周长的非圆截面与圆截面，不误作同面积。
\begin{tikzpicture}
  \draw[body,fill=blue!5] (0,0) ellipse[x radius=1.4,y radius=.45];
  \node at (0,0) {$S$};\node[below] at (0,-.8) {周长 $P$，非圆截面};
  \draw[body,fill=green!7] (4,0) circle(.9870588543);
  \node at (4,0) {$P^2/(4\pi)$};\node[below] at (4,-1.25) {同周长圆截面};
  \draw[vec] (1.7,0)--(2.7,0);
  \node[above] at (2.2,1.25) {$S\le P^2/(4\pi)$};
\end{tikzpicture}
\end{document}
